\section{Characteristic lengths}\label{df:sec:scales}

The rare-fast estimate concerns crossings with a negative defect.
We next identify a length at which crossings are rare even with a
positive allowance of order the strip width. A finite-region change
of measure supplies an upper bound on that length. Its definition
then guarantees enough shorter crossings to construct the endpoint
extensions in the next section.

Small multiplicative changes of many independent weights, with a
controlled quadratic probability cost, also appear in fluctuation
arguments; see Chatterjee \cite{Chatterjee2019}. Here the exact Gamma
likelihood calculation is valid for every positive shape parameter.

\subsection{A first crossing threshold}

Set
\[
 K=100\bigl(1+\abs{a_-}+\abs{a_+}\bigr).
\]
For real $s\geq1$ and integer $n\geq1$, let $E(s,n)$ be the event that
there is a finite path $P$ from the column $x=0$ to the column $x=n$,
with every vertex in
\[
 \{-n\leq x\leq2n,\ \abs{r}\leq8s\},
\]
and with
\[
 \tau(P)<m_0n+Ks.
\]
Endpoints on the two columns are unrestricted within this region. As
throughout, the event tests only the internal edges of the indicated
region. The notation $E(s,n)$ will also be used with an explicitly
specified lattice translation.

\begin{lemma}[Rarity at cubic length]\label{df:gr:tilt}
For $s=2^k$ and $k\to\infty$,
\[
 \PP\bigl(E(s,s^3)\bigr)\longrightarrow0.
\]
\end{lemma}
\begin{proof}
\emph{Turning an allowed crossing into a fast crossing.}
Put $n=s^3$. Choose a fixed $\eta>0$ such that
\[
 \eta m_0>K+16\abs{a_+}+1,
\]
and set $\delta=\eta s/n$ and $\alpha=1-\delta$. Take $s$ sufficiently
large that $0<\delta\le1/2$. Multiply every internal edge weight of the
region defining $E(s,n)$ by $\alpha$, leaving the other weights
unchanged. If $P$ witnesses $E(s,n)$ before this operation, its
transverse displacement obeys $\abs{r(\Delta P)}\le16s$, and in the
modified environment
\[
 \begin{split}
 H_+(P)
 &<\alpha(m_0n+Ks)-m_0n-a_+r(\Delta P)\\
 &\le \bigl(K-\eta m_0+16\abs{a_+}\bigr)s
          -\eta K\frac{s^2}{n}<-s.
 \end{split}
\]
This computation includes every traversal if $P$ repeats an edge.
Let $F(s,n)$ be the crossing event $H_+(P)<-s$ in the same region, and
let $Q$ be the law of the modified environment. The deterministic
implication gives
\begin{equation}\label{df:eq:gr-tilt-implication}
 \PP(E(s,n))\le Q(F(s,n)).
\end{equation}

\emph{The cost of the change of measure.}
Under $Q$, a modified edge has distribution
$\operatorname{Gamma}(\kappa,\lambda/\alpha)$. On $t>0$ its likelihood
ratio with respect to the original edge law is
\[
 L_\delta(t)
 =\frac{p_{\kappa,\lambda/\alpha}(t)}{p_{\kappa,\lambda}(t)}
 =\alpha^{-\kappa}
    \exp\{-\lambda(\alpha^{-1}-1)t\}.
\]
Both densities are positive on $(0,\infty)$, including when
$0<\kappa<1$. Direct integration gives
\[
 \begin{split}
 \E[L_\delta(\tau_e)^2]
 &=\alpha^{-2\kappa}\frac{\lambda^\kappa}{\Gamma(\kappa)}
   \int_0^\infty t^{\kappa-1}
     e^{-\lambda(2/\alpha-1)t}\,\mathrm d t\\
 &=\alpha^{-2\kappa}(2/\alpha-1)^{-\kappa}
   =\bigl(\alpha(2-\alpha)\bigr)^{-\kappa}
   =(1-\delta^2)^{-\kappa}.
 \end{split}
\]
The region has at most $(3n+1)(16s+2)$ vertices and at most twice that
many internal unoriented edges. Thus the number $M$ of modified
coordinates is at most $144ns$. The likelihood ratio is a product
over these distinct coordinates, regardless of path multiplicity.
By independence,
\[
 \begin{split}
 \E\left[\left(\frac{\mathrm d Q}{\mathrm d\PP}\right)^2\right]
 &=(1-\delta^2)^{-\kappa M}
 \le \exp(2\kappa M\delta^2)\\
 &\le \exp\left(288\kappa\eta^2\frac{s^3}{n}\right)
 =\exp(288\kappa\eta^2).
 \end{split}
\]
\emph{Transfer back to the original law.}
Lemma~\ref{df:cr:rare-fast}, applied under the original law with $A=1$,
$R=8$, and $\varepsilon=1$, gives $\PP(F(s,n))\to0$ because $n/s=s^2\to\infty$.
Cauchy--Schwarz therefore gives
\[
 \PP(E(s,n))\le Q(F(s,n))
 \le e^{144\kappa\eta^2}\PP(F(s,n))^{1/2}\longrightarrow0.
\]
All supports and the norm in $F(s,n)$ are those of the original fixed
law; no shape theorem for the modified environment is used.
\end{proof}

Fix any probability threshold $q_0\in(0,1)$ for this section and the
next. It may be arbitrarily small; Section~\ref{df:sec:witnesses}
will specify its value. Write $s_k=2^k$.

\begin{lemma}[Characteristic lengths]\label{df:gr:characteristic-lengths}
For every fixed $C>0$,
\begin{equation}\label{df:eq:gr-short-crossings}
 \PP(E(s,n))\longrightarrow1
 \qquad\text{as }s\to\infty,\quad 1\leq n\leq Cs,\quad n\in\N.
\end{equation}
There is therefore an integer $k_0\geq0$ such that, for every $k\geq k_0$,
the first dyadic integer $L_k\in\{1,2,4,\ldots\}$ satisfying
\begin{equation}\label{df:eq:gr-characteristic-threshold}
 \PP(E(s_k,L_k))\leq q_0
\end{equation}
exists. These lengths satisfy
\begin{equation}\label{df:eq:gr-characteristic-bounds}
 L_k\leq2^{3k},\qquad \frac{L_k}{2^k}\longrightarrow\infty.
\end{equation}
Moreover, every dyadic integer $l<L_k$ satisfies
$\PP(E(s_k,l))>q_0$.
\end{lemma}

\begin{proof}
\emph{Short crossings are likely.}
First suppose $n\to\infty$ with $n\leq Cs$, and put
$q=(n,\lfloor bn\rfloor)$. Then
\[
 \mu(q)\leq m_0n+\mu(e_2),\qquad -1<r(q)\leq0.
\]
Apply Lemma~\ref{df:cr:local-travel} at scale $n$, with endpoints $0,q$,
with cost error $\eta n$ where $\eta=K/(4C)$, and with tube radius
$\delta_0n$, where
\[
 0<\delta_0\leq\min\{1/4,1/(2C)\}.
\]
The endpoints have Euclidean distance at most $3n$. Since
$\abs{b}\leq1$, changing a point by Euclidean distance $t$ changes its
$r$-coordinate by at most $\sqrt2t$. The resulting tube is inside
$-n\leq x\leq2n$, and inside $\abs{r}\leq8s$ for all sufficiently large
$s$: its transverse extent is at most $1+\sqrt2\delta_0n$. With
probability tending to one its path cost is at most
\[
 m_0n+\mu(e_2)+\eta n
 \leq m_0n+\mu(e_2)+Ks/4
 <m_0n+Ks.
\]
This proves \eqref{df:eq:gr-short-crossings} when $n\to\infty$.

If $n$ remains bounded, direct lattice paths from $0$ to
$(n,\lfloor bn\rfloor)$ use a bounded number of edges and remain in a
bounded transverse region with $0\leq x\leq n$. Their costs are bounded
in probability, whereas $Ks\to\infty$. This proves the same assertion
for bounded $n$. The two cases give the joint limit in
\eqref{df:eq:gr-short-crossings} by the subsequence criterion.

\emph{The first dyadic threshold.}
Lemma~\ref{df:gr:tilt} shows that the dyadic integer $2^{3k}$ satisfies
\eqref{df:eq:gr-characteristic-threshold} for all sufficiently large $k$.
Thus $L_k$ exists and has the stated upper bound. If $L_k/2^k$ failed to
tend to infinity, a subsequence would satisfy $L_k\leq C2^k$ for some
fixed $C$, contradicting \eqref{df:eq:gr-short-crossings} and
\eqref{df:eq:gr-characteristic-threshold}. The last assertion follows from
the definition of the first dyadic length; it uses no monotonicity in
$n$ of the events $E(s,n)$.
\end{proof}

\subsection{Scales compatible with endpoint extensions}

The sequence $L_k$ need not be monotone. The endpoint construction
only requires a lower bound for earlier lengths relative to the one
currently in use. Record minima give this on an infinite subsequence.

\begin{lemma}[Record scales]\label{df:gr:record-scales}
There is an infinite set $\mathcal K\subset\{k_0,k_0+1,\ldots\}$ such that,
for $k\in\mathcal K$ and every integer $0\leq j\leq k-k_0$,
\begin{equation}\label{df:eq:gr-record-comparison}
 L_{k-j}\geq2^{-4j}L_k.
\end{equation}
\end{lemma}

\begin{proof}
The positive sequence $L_k2^{-4k}$ is at most $2^{-k}$, so it tends to
zero and has infinitely many strict record minima. Take these indices
for $\mathcal K$. At a record index $k$, comparison with $k-j$ gives
\[
 L_k2^{-4k}\leq L_{k-j}2^{-4(k-j)},
\]
which is \eqref{df:eq:gr-record-comparison}.
\end{proof}

Henceforth $k$ tends to infinity through $\mathcal K$, and we abbreviate
\[
 s=2^k,\qquad L=L_k=2^m.
\]
In particular $m\leq3k$ and $m-k\to\infty$.
